% ============================================================
% 経営情報 2026　学習ガイド　共通プリアンブル
%   study-guide.tex から読み込む。体裁は materials/ai-review/preamble.tex に合わせている。
% ============================================================
\documentclass[paper=a4paper,fontsize=10pt,line_length=47zw,number_of_lines=43,baselineskip=1.6zw]{jlreq}

\usepackage[haranoaji]{luatexja-preset}
\usepackage{xcolor}
\usepackage[most]{tcolorbox}
\usepackage{tabularray}
\usepackage{enumitem}
\usepackage{needspace}
\usepackage{fancyhdr}
\usepackage[hidelinks,bookmarksdepth=2,bookmarksnumbered=false]{hyperref}
\hypersetup{pdftitle={経営情報 2026 学習ガイド},pdfauthor={河合勝彦}}

% ------------------------------------------------------------
% 色（サイトの assets/css/style.css に合わせた基調色と、モジュールごとの色）
% ------------------------------------------------------------
\definecolor{ink}{HTML}{1D2330}
\definecolor{inktwo}{HTML}{4A5261}
\definecolor{inkthree}{HTML}{757D8A}
\definecolor{rulegray}{HTML}{B9BEB4}
\definecolor{papertwo}{HTML}{ECEEE8}
\definecolor{bad}{HTML}{B3261E}    \definecolor{badtint}{HTML}{F9DEDA}
\definecolor{mone}{HTML}{1E4A4A}   \definecolor{monetint}{HTML}{E3EFEA}
\definecolor{mtwo}{HTML}{B5651D}   \definecolor{mtwotint}{HTML}{F7E7D5}
\definecolor{mthree}{HTML}{2F5D8A} \definecolor{mthreetint}{HTML}{E0E9F3}
\definecolor{mfour}{HTML}{6B3A6B}  \definecolor{mfourtint}{HTML}{EEE2EE}
\colorlet{phase}{mone}
\colorlet{phasetint}{monetint}

% ------------------------------------------------------------
% ページスタイル
% ------------------------------------------------------------
\newcommand*\sessionmark{}
\pagestyle{fancy}
\fancyhf{}
\renewcommand{\headrulewidth}{0pt}
\fancyfoot[L]{\sffamily\scriptsize\color{inkthree}経営情報 2026　学習ガイド}
\fancyfoot[R]{\sffamily\scriptsize\color{inkthree}\sessionmark\quad\thepage}
\fancypagestyle{cover}{\fancyhf{}}

\setlist{nosep,leftmargin=1.5em}
\setcounter{tocdepth}{1}

% ------------------------------------------------------------
% 各回の見出し　\session{回}{日付}{時限}{モジュール 1-4}{題目}
% ------------------------------------------------------------
\newcommand*\modulename{}
\newcommand\session[5]{%
  \clearpage
  \ifcase#4\or
    \colorlet{phase}{mone}\colorlet{phasetint}{monetint}\renewcommand*\modulename{Module 01　経営と情報の基礎}\or
    \colorlet{phase}{mtwo}\colorlet{phasetint}{mtwotint}\renewcommand*\modulename{Module 02　企業と情報システム}\or
    \colorlet{phase}{mthree}\colorlet{phasetint}{mthreetint}\renewcommand*\modulename{Module 03　データと意思決定}\or
    \colorlet{phase}{mfour}\colorlet{phasetint}{mfourtint}\renewcommand*\modulename{Module 04　デジタル時代の経営}\fi
  \renewcommand*\sessionmark{第#1回　#5}%
  \phantomsection\addcontentsline{toc}{section}{\texorpdfstring{第#1回　#5\hfill{\protect\small #2　#3}}{第#1回　#5}}%
  \begin{tcolorbox}[enhanced,colback=phase,colframe=phase,arc=2mm,boxrule=0pt,
      left=4mm,right=4mm,top=2.5mm,bottom=2.5mm,fontupper=\color{white}\sffamily]
    {\small 第#1回　#2　#3　\textbar　\modulename}\par
    \vspace{0.6mm}{\LARGE\bfseries #5}
  \end{tcolorbox}
  \vspace{1mm}%
}

% 序章などの見出し（目次には第1階層で載せる）
\newcommand\intro[1]{%
  \clearpage
  \renewcommand*\sessionmark{#1}%
  \phantomsection\addcontentsline{toc}{section}{#1}%
  {\noindent\sffamily\LARGE\bfseries\color{phase}#1}\par
  \vspace{1mm}\noindent{\color{phase}\rule{\linewidth}{1pt}}\par\vspace{3mm}}

% 各回の中の節（解説・演習課題・模範解答・クイズ解説）
\newcommand\guidesec[1]{%
  \par\needspace{32mm}\vspace{4mm}%
  \phantomsection\addcontentsline{toc}{subsection}{#1}%
  {\noindent\sffamily\large\bfseries\color{phase}#1}\par
  \vspace{0.6mm}\noindent{\color{phase}\rule{\linewidth}{0.6pt}}\par\vspace{2mm}}

% 小見出し
\newcommand\topic[1]{%
  \par\needspace{18mm}\vspace{2.5mm}%
  {\noindent\sffamily\bfseries\color{ink}#1}\par\vspace{0.6mm}}

% ねらい
\newcommand\goals[1]{%
  \begin{tcolorbox}[enhanced,colback=phasetint,colframe=phasetint,arc=1.5mm,boxrule=0pt,
      left=3mm,right=3mm,top=1.5mm,bottom=1.5mm,fontupper=\small]
    {\sffamily\bfseries\color{phase}この回のねらい}\par
    \begin{itemize}[label=\textcolor{phase}{・},leftmargin=1.2em]#1\end{itemize}
  \end{tcolorbox}}

% 演習の見出し　\exercise{番号}{題}
\newcommand\exercise[2]{%
  \par\needspace{30mm}\vspace{3mm}%
  {\noindent\sffamily\bfseries\color{phase}■ 演習#1}\enspace{\sffamily\bfseries #2}\par\vspace{1mm}}

% 指示文
\newcommand\instr[1]{{\small\color{inktwo}#1}\par\vspace{1mm}}

% プロンプト例
\newtcolorbox{promptbox}[1][]{enhanced,colback=papertwo,colframe=papertwo,boxrule=0pt,
  arc=1mm,leftrule=1.2mm,colbacktitle=papertwo,
  borderline west={1.2mm}{0pt}{phase},left=3mm,right=2mm,top=1mm,bottom=1mm,
  fontupper=\sffamily\footnotesize\color{inktwo},before skip=1mm,after skip=2mm,
  title={\sffamily\scriptsize\bfseries\color{phase}プロンプト例　#1},
  attach title to upper={\par\vspace{0.5mm}}}

% 模範解答　\begin{answer}{演習1}
\newtcolorbox{answer}[1]{enhanced,breakable,colback=white,colframe=phase,boxrule=0.6pt,arc=1.5mm,
  left=3mm,right=3mm,top=1.5mm,bottom=1.5mm,fontupper=\small,
  before skip=2mm,after skip=3mm,
  title={\sffamily\bfseries 模範解答　#1},colbacktitle=phase,coltitle=white,
  fonttitle=\small,toptitle=0.8mm,bottomtitle=0.8mm}

% 注意
\newtcolorbox{caution}[1]{enhanced,colback=badtint,colframe=badtint,boxrule=0pt,arc=1mm,
  left=3mm,right=3mm,top=1mm,bottom=1mm,fontupper=\footnotesize,
  before upper={{\sffamily\bfseries\color{bad}#1}\par}}

% ポイント
\newtcolorbox{point}[1]{enhanced,colback=phasetint,colframe=phasetint,boxrule=0pt,arc=1mm,
  left=3mm,right=3mm,top=1mm,bottom=1mm,fontupper=\small,
  before upper={{\sffamily\bfseries\color{phase}#1}\par}}

% チェック項目
\newenvironment{checks}{\begin{itemize}[label={□},leftmargin=1.5em,itemsep=0.8mm]\small}{\end{itemize}}

% 表の共通スタイル
\NewTblrEnviron{gtab}
\SetTblrInner[gtab]{hlines={0.4pt,rulegray},vlines={0.4pt,rulegray},
  cells={font=\small},rowsep=2pt,colsep=4pt}
\NewTblrEnviron{gtabh}
\SetTblrInner[gtabh]{hlines={0.4pt,rulegray},vlines={0.4pt,rulegray},
  cells={font=\small},rowsep=2pt,colsep=4pt,
  row{1}={bg=phasetint,font=\sffamily\footnotesize\bfseries}}

% 表の前後の余白（表は \par\noindent で段落の先頭に置く）
\AddToHook{env/gtab/after}{\par\vspace{1.2mm}}
\AddToHook{env/gtabh/after}{\par\vspace{1.2mm}}
\AddToHook{env/keywords/after}{\par\vspace{1.2mm}}

% キーワード表　\begin{keywords}{} 用語 & 意味 \\ ... \end{keywords}（tabularray の環境なので空の引数 {} が要る）
\NewTblrEnviron{keywords}
\SetTblrInner[keywords]{colspec={Q[l,wd=34mm] X[l]},hlines={0.4pt,rulegray},vlines={0.4pt,rulegray},
  cells={font=\small},rowsep=2pt,colsep=4pt,column{1}={font=\sffamily\small\bfseries}}

% ------------------------------------------------------------
% 出席確認クイズ　\quizq{番号}{問題文}{選択肢a}{選択肢b}{選択肢c}{選択肢d}{解説}
%   正解の選択肢は \correct{...} で囲む（quiz/NN.tex は extract スクリプトで生成）。
% ------------------------------------------------------------
\newcommand\correct[1]{\textbf{#1}\enspace{\sffamily\bfseries\color{phase}［正解］}}
\newcommand\quizq[7]{%
  \par\needspace{34mm}\vspace{2.5mm}%
  {\noindent\sffamily\bfseries\color{phase}Q#1.}\enspace #2\par
  \begin{enumerate}[label=(\alph*),leftmargin=2.4em,labelsep=0.5em,itemsep=0.2mm,topsep=0.6mm]
    \item #3 \item #4 \item #5 \item #6
  \end{enumerate}
  \begin{tcolorbox}[enhanced,breakable,colback=phasetint,colframe=phasetint,boxrule=0pt,arc=1mm,
      left=2.5mm,right=2.5mm,top=1mm,bottom=1mm,fontupper=\small,before skip=1mm,after skip=2mm]
    {\sffamily\bfseries\color{phase}解説}\enspace #7
  \end{tcolorbox}}
